\begin{statement}{A.25}{Proved}{The stable genus-two normalizer at every odd prime}
Let
\[
 D:\ z^2=(t^2+4)(t^3+3t),\qquad
 Y=D\setminus\{Q_\infty\},
\]
and let $\ell$ be odd. Every nontrivial closed procyclic
subgroup of $\pi_1(Y_{\overline K})^{(\ell)}$ whose
arithmetic normalizer has open image in $G_K$ is contained
in a conjugate of cusp inertia.
Together with A.14 at $\ell=2$, this completes the
pro-$\ell$ normalizer classification for every first cover
and every displayed genus-two quotient in Task A.4.

The extra detector is the degree-two separating monodromy
at a place above $2$, acting on the proper genus-two group.
It is an actual arithmetic derivation and does not require
a second cusp-compatible elliptic map or a lifting through
the odd-prime quotient $P_{C,\ell}\twoheadrightarrow P_{D,\ell}$.
\end{statement}
\begin{proof}
\paragraph{The degree-three map.}
The map to the original elliptic curve is
\[
 f:D\longrightarrow E,\qquad
 x_E=-\frac{i}{2}(t^3+3t),\qquad
 y_E=-\frac{1+i}{4}z(t^2+1).
                                                        \tag{A.25.1}
\]
Its coordinates satisfy $y_E^2=x_E^3-x_E$.
The pole order of $x_E$ at the unique point at infinity
is six, so $\deg f=3$ and
\[
                         f^*(O)=3Q_\infty.              \tag{A.25.2}
\]
Riemann--Hurwitz then shows that the indicated totally
ramified cusp is all the ramification. In particular
$f:Y\to E\setminus O$ is finite etale. Formula (A.25.1)
also follows by descending the $C_3$ map in A.4:
$w^3-w^{-3}=t^3+3t$ and
$w^3+w^{-3}=(w+w^{-1})(t^2+1)$.

\paragraph{The stable reduction above two.}
Work over a finite extension of $K_2$ on which $E$ has
good reduction, $D$ has stable reduction, and all
components are defined. Such an extension exists.
The potential good reduction of $E$ follows from its
integral $j$-invariant; the exact criterion is
\cite[Proposition VII.5.5]{fghSilvermanAEC}.
Its geometric special fibre $E_0$ has $j=0$ in
characteristic two and is supersingular. For example,
over an algebraic closure its equation is
$y^2+y=x^3$, whose only geometric two-torsion point is
the origin. Hence there is no etale degree-two isogeny
onto $E_0$.

By A.4, $J_D\sim E^2$ after a finite extension.
Good reduction of an abelian variety is invariant
under isogeny, as follows from the
Neron--Ogg--Shafarevich criterion. Thus $J_D$ has good
reduction here. The dual graph of the stable fibre of
$D$ is a tree: for a nodal curve the torus in its
generalized Jacobian has rank the first Betti number
of the dual graph. Consequently the stable fibre is
either smooth of genus two or two smooth elliptic
curves meeting in one separating node. In genus two
there are no further stable compact-type possibilities:
a rational vertex of valence at least three would
require at least three positive-genus leaves.

The map $f$ extends to this stable model with values in
the good elliptic model $\mathcal E$. Here is the
extension argument, including the nodal case.
Resolve the rational map on a regular semistable model
by blowing up closed points. Every exceptional
geometric curve is rational and maps constantly to
the elliptic special fibre. The same is true of every
rational component contracted to obtain the stable
model. The map is therefore constant on the connected
fibres of the contraction. Its graph image over the
stable model is proper and quasi-finite, hence finite,
and birational. The stable total space is normal
(locally at a node it is $xy=\varpi^e$), so this finite
birational graph is the stable model itself. This
proves descent of the morphism.

Suppose first that the stable fibre $D_0$ were smooth.
The pullback of $\mathcal O_{\mathcal E}(O)$ has degree
three on every fibre, so $f_0:D_0\to E_0$ has degree
three and is nonconstant. There is no vertical
component in the Cartier divisor $f^*(O)$.
Equation (A.25.2) therefore extends as the relative
divisor identity $f^*(O)=3\overline{Q_\infty}$, and
\[
                         f_0^*(O)=3(Q_\infty)_0.
\]
The degree-three map is separable in characteristic
two. Its ramification index three at this point is
tame and contributes two to the different.
Riemann--Hurwitz gives total different degree two,
so there is no other ramification.

The Galois closure of this separable cubic map has
group $C_3$ or $S_3$. In the cyclic case a Kummer
description would require a rational function on
$E_0$ whose divisor has exactly one valuation
nonzero modulo three, contradicting degree zero.
In the $S_3$ case inertia at the sole branch point
is $C_3$, and the quotient by $A_3$ is a connected
etale double cover of $E_0$. It is a degree-two
isogeny after a choice of origin, contradicting
supersingularity. Thus $D_0$ is not smooth.

We have proved that
\[
                  D_0=E_1\cup_N E_2                    \tag{A.25.3}
\]
with both components elliptic. The sum of the degrees
of the two restrictions $f_i:E_i\to E_0$ is three,
by the same fibrewise degree calculation. Neither
restriction is constant. Indeed, suppose the degrees
were $(3,0)$. The nonconstant restriction is a
separable degree-three isogeny of elliptic curves,
up to translation, hence is etale and has three
distinct geometric preimages of $O$.
On the open part $E_1\setminus\{N\}$, no vertical
component of the Cartier divisor $f^*(O)$ is present:
the nonconstant component itself is not a component
of this divisor, and the only other component meets
it at $N$. Its horizontal part is exactly
$3\overline{Q_\infty}$. Thus the support of
$f_1^{-1}(O)$ outside $N$ has at most one point.
An etale degree-three isogeny has at least two of
its three preimages outside $N$, a contradiction.
This argument also covers specialization of the
cusp to the node or to the contracted component.
Interchanging the indices excludes $(0,3)$.
After relabelling, therefore,
\[
                         \deg f_1=1,\qquad \deg f_2=2.
                                                        \tag{A.25.4}
\]
Both restrictions are isogenies after translation.
Their maps on rational $\ell$-adic homology are
isomorphisms for every odd $\ell$, including when
the degree-two map in characteristic two is
inseparable.

\paragraph{The two nonzero components of a candidate line.}
Suppose $A$ is a noncuspidal arithmetically normalized
procyclic subgroup of the pro-$\ell$ group of $Y$.
By A.11 a generator has a nonzero proper homology
class $v$. The map (A.25.1), together with the
elliptic normalizer criterion A.12, gives
$f_*v=0$: a nonzero image would produce a
noncuspidal normalized procyclic subgroup on
$E\setminus O$.

Good reduction of the Jacobian and compact-type
specialization identify
\[
 H_1(D_{\overline K},\Q_\ell)=V_1\oplus V_2,\qquad
 V_i=H_1(E_i,\Q_\ell).
                                                        \tag{A.25.5}
\]
This identification is compatible with $f_*$, whose
restrictions are $(f_i)_*$. By (A.25.4) both are
isomorphisms. Hence
\[
                  v=v_1+v_2,\qquad v_1\ne0,\quad v_2\ne0.
                                                        \tag{A.25.6}
\]
For example, if $v_1=0$, the equality
$(f_1)_*v_1+(f_2)_*v_2=0$ and injectivity of
$(f_2)_*$ force $v_2=0$.

\paragraph{The actual monodromy and its Lie obstruction.}
Pass now to the proper pro-$\ell$ group of $D$.
The image of $A$ is still a nontrivial procyclic
group, since its homology class is $v\ne0$, and
its arithmetic normalizer still has open image.
This passage does not assert an injective comparison
from a prime-to-$\ell$ covering group.

The exact monodromy input is the non-abelian
Picard--Lefschetz theorem
\cite[Definition 3.1.5 and Theorem 3.1.6]{fghBettsDogra}.
For a regular semistable model it identifies the
action of tame inertia on the rational unipotent
fundamental group with edge twists of its reduction
graph. Resolving the node (A.25.3) merely subdivides
its one edge; its total length is a positive number
$m$. Choose a base point specializing to the second
elliptic component, after another finite extension
if needed. The twist fixes its vertex group and
conjugates the first vertex group by the $m$-th
power of the separating boundary. Its first
Johnson degree is two. With
\[
 \omega_1=[a,b],\qquad \omega_2=[c,d],\qquad
 \operatorname{gr}\mathfrak p_D
       =\mathbb L(a,b,c,d)/(\omega_1+\omega_2),
                                                        \tag{A.25.7}
\]
the degree-two derivation, up to the harmless common
nonzero tame scalar and an inner derivation, is
\[
        T(x)=m[\omega_1,x]\ (x\in V_1),\qquad
        T(x)=0\ (x\in V_2).                              \tag{A.25.8}
\]
The graph presentation in Definition 3.1.5 gives
(A.25.8) directly; it does not infer a Galois
derivation from an ambient formal possibility.

We record the necessary normalizer condition
carefully. An open arithmetic image contains a
nonzero power of a tame inertia element after all
the finite extensions made above. Inertia acts
trivially on the proper homology, since the
Jacobian has good reduction. Thus its conjugation
character on the proposed procyclic subgroup is
one. In the rational unipotent group, its standard
action must consequently carry the image of a
generator $u$ to a conjugate of $u$.
The action in (A.25.8) has zero Johnson degree one.
Comparing degree two of this conjugacy equation
gives $[q_1,v]=0$ in
$\Lambda^2(V_1\oplus V_2)/\Q_\ell(\omega_1+\omega_2)$,
where $q_1$ is the linear term of a conjugator.
A decomposable bivector cannot be a nonzero
multiple of the nondegenerate rank-four symplectic
form, so $q_1$ is a scalar multiple of $v$.
Multiplying the rational conjugator by the inverse
of the corresponding rational power of $u$ removes
that term. This is only a necessary equation in a
rational quotient, not a claimed integral lift.
Comparison in degree three now yields
\[
            T(v)\in[v,\operatorname{gr}_2\mathfrak p_D].
                                                        \tag{A.25.9}
\]
Changing the initial representative of the outer
monodromy adds an inner derivation and does not
change this necessary condition.

Choose symplectic bases with $a=v_1$ and $c=v_2$;
this is possible by (A.25.6). There is a Lie map
from the algebra in (A.25.7) to the free Lie
algebra on $x,y$ defined by
\[
                 a\mapsto x,\quad c\mapsto-x,\quad
                 b\mapsto y,\quad d\mapsto y.
                                                        \tag{A.25.10}
\]
It kills both $v=a+c$ and the relation
$[a,b]+[c,d]$. But it carries $T(v)$ to
\[
                          m[[x,y],x]\ne0.
                                                        \tag{A.25.11}
\]
It therefore kills the right side of (A.25.9)
and not the left, a contradiction. This excludes
the proposed subgroup at every odd prime.

All nontrivial closed subgroups of cusp inertia
do have open arithmetic normalizer after the
finite extension making that cusp rational.
This proves the stated classification. The
independent exact certificate
\texttt{code/a\_D3\_monodromy.py} verifies (A.25.1),
the origin-fibre pole orders, and the nonzero
free-Lie witness (A.25.10)--(A.25.11).
\end{proof}
